\section{Evidence Catalogue and Consistency Audit}
\label{app:evidence-audit}

\graphicspath{{../evidence/}{evidence/}}

This appendix connects the analytical report to the separately preserved
public evidence package. It has three purposes: to identify the evidence
snapshot used during report preparation, to record the result of an independent
manifest and consistency check, and to provide a visual index of every public
screenshot without treating the report itself as the raw-evidence store.

The full-resolution PNG captures, command output, source, configuration,
generated objects, binaries, and phase manifests remain under
\path{evidence/}. This revision renders the original screenshot files directly
from that tree. The filename shown below each capture is therefore also the
repository-relative source path. The figures provide visual coverage, while the
phase manifests remain authoritative for exact bytes and integrity checks.

% One declared visual sheet now occupies exactly one physical page.  Version 1
% used ordinary float pages, which allowed LaTeX to combine two short sheets and
% occasionally pushed their contents into the header area.
\newcommand{\evidencepath}[1]{%
    \par\smallskip
    \begingroup
        \centering
        \scriptsize\ttfamily\color{MoriMuted}%
        \expandafter\path\expandafter{\detokenize{#1}}\par
    \endgroup
}

\newcommand{\evidencepair}[3]{%
    \begin{figure}[H]
        \centering
        \includegraphics[
            width=\textwidth,
            height=0.35\textheight,
            keepaspectratio
        ]{#1}
        \evidencepath{#1}
        \vspace{1.1em}
        \includegraphics[
            width=\textwidth,
            height=0.35\textheight,
            keepaspectratio
        ]{#2}
        \evidencepath{#2}
        \caption{#3}
    \end{figure}
    \clearpage
}

\newcommand{\evidencesidebyside}[3]{%
    \begin{figure}[H]
        \centering
        \begin{minipage}[t]{0.48\textwidth}
            \centering
            \includegraphics[
                width=\linewidth,
                height=0.68\textheight,
                keepaspectratio
            ]{#1}
            \evidencepath{#1}
        \end{minipage}%
        \hfill
        \begin{minipage}[t]{0.48\textwidth}
            \centering
            \includegraphics[
                width=\linewidth,
                height=0.68\textheight,
                keepaspectratio
            ]{#2}
            \evidencepath{#2}
        \end{minipage}
        \caption{#3}
    \end{figure}
    \clearpage
}

\newcommand{\evidencesingle}[2]{%
    \begin{figure}[H]
        \centering
        \includegraphics[
            width=\textwidth,
            height=0.76\textheight,
            keepaspectratio
        ]{#1}
        \evidencepath{#1}
        \caption{#2}
    \end{figure}
    \clearpage
}

\subsection{Evidence Snapshot and Authority}

The audit used the intact repository snapshot at commit:

\begin{lstlisting}
d78fce4f0aa9554a5f9f19503b75056f75209619
\end{lstlisting}

That snapshot contains the complete public evidence tree through MORI v2.7 and
Phase~05. Two more recent
uploaded ZIP copies lacked their end-of-central-directory records and ended
inside compressed entries. A controlled local-header recovery found 185 and
147 entries respectively, but each final recovered entry failed decompression.
They were therefore used only for comparison and were not promoted over the
intact repository snapshot. A third, older ZIP was structurally valid but ended
at an earlier report and evidence state.

This selection rule avoids two opposite errors: accepting a newer filename as
proof of completeness, or silently falling back to an older package without
checking whether later MORI and vendor-validation evidence exists elsewhere.
The ZIP digests and recovery results are retained in the accompanying audit
supplement.

\subsection{Public-Manifest Verification}
\label{app:evidence-manifest-audit}

The seven selected public manifests contain 220 unique entries. The Phase~04
manifest already repeats the 108 entries from the nested MORI v2 manifest, so
the nested list is verified but not added to the total a second time. The audit
result is shown in \cref{tab:evidence-manifest-audit}.

\begin{table}[H]
    \centering
    \caption{Independent verification of the authoritative public manifests.}
    \label{tab:evidence-manifest-audit}
    \small
    \begin{tabularx}{\textwidth}{@{}>{\bfseries\raggedright\arraybackslash}p{0.28\textwidth}>{\centering\arraybackslash}p{0.12\textwidth}>{\centering\arraybackslash}p{0.12\textwidth}Y@{}}
        \toprule
        Evidence set & Listed & Verified & Recorded exception \\
        \midrule
        Phase 00 baseline & 20 & 20 & None. \\
        Phase 01 validation & 14 & 14 & None. \\
        Phase 02 reproduction & 12 & 11 & The listed \code{README.md} digest is stale. \\
        Phase 03 detection & 22 & 21 & The listed \code{README.md} digest is stale. \\
        Phase 03 YARA comparison & 8 & 7 & The listed \code{README.md} digest is stale. \\
        Phase 04 custom mitigation & 126 & 126 & None. \\
        Phase 05 patched validation & 18 & 18 & None. \\
        \midrule
        Total & 220 & 217 & Three documentation-only mismatches. \\
        \bottomrule
    \end{tabularx}
\end{table}

Every non-README entry in the three affected manifests verified. Those three
README files are consequently used as explanatory documentation, not as
integrity-verified experimental objects. The report does not regenerate the
original manifests and pretend that the earlier frozen lists never existed;
the exact expected and observed README digests are preserved in the audit
supplement.

A second documentation mismatch concerns inventory rather than content. The
inner MORI v2 README still states 45 artifacts and 40 screenshots, for an
85-object package. The current verified MORI v2 manifest instead contains 49
implementation or provenance artifacts and 59 screenshots, for 108 objects.
The manifest, directory contents, and Phase~04 index agree on 108, so that is
the authority used in Chapters~7 and~10.

\subsection{VM 270 Hypervisor Supplement}

The operator-supplied Proxmox transcript for VM~270 was reconstructed from the
terminal output provided during report preparation. Rehashing the exact byte
sequence produced:

\begin{lstlisting}
698bfb469060ea9bc0b9d5afdc8636c7c551db9b3f388181a4353f710e501073
\end{lstlisting}

This equals the digest displayed by the original \code{sha256sum} command. The
capture records VM~270 as stopped on 13 August 2026 at 21:53:07+02:00, names
\code{clean-baseline-00} as its parent, and places the current state below that
snapshot. The exact transcript contains MAC, storage-volume, SMBIOS, and
VM-generation identifiers, so it remains in the private audit set. The
sanitised publication copy retains the evidential fields used in Chapter~3 and
has SHA-256:

\begin{lstlisting}
706cbe7fa1a2700ff00d7b7ac5394ddeed8811e5ba5397fa8b71d830e9c3089c
\end{lstlisting}

\subsection{Phase 00: Clean Reference Baseline}
\label{app:evidence-phase00}

\begin{evidencebox}
\textbf{Public path:} \path{evidence/00-baseline/}

\textbf{Manifest:} \path{evidence/00-baseline/hashes/PUBLIC-SHA256SUMS}

\textbf{Coverage:} 20 verified command-output artifacts covering guest
identity, kernel and package state, networking, listeners, storage, and the
recorded kernel-image digest. The phase has no public screenshot directory.

\textbf{Boundary:} the baseline fixes the guest's documented starting state;
it does not independently establish the semantic correctness of each command,
the contents of the unexported VM disk, or Copy Fail exploitability.
\end{evidencebox}

The private and public baseline manifests intentionally describe different
byte sets where publication redaction changed guest identifiers. The public
manifest is repeatable from the publication package; the historical private
verification depends on the separately retained originals. The VM~270
hypervisor supplement adds snapshot existence and ancestry without turning a
snapshot label into proof of every guest-state claim.

\subsection{Phase 01: Reference-System Validation}
\label{app:evidence-phase01}

\begin{evidencebox}
\textbf{Public path:} \path{evidence/01-vulnerability-validation/}

\textbf{Manifest:}
\path{evidence/01-vulnerability-validation/PUBLIC-SHA256SUMS}

\textbf{Coverage:} seven verified text artifacts and seven verified
screenshots covering the running kernel, package versions, module metadata,
kernel configuration, effective module policy, dry-run resolution, and the
negative loaded-module result.

\textbf{Boundary:} no proof of concept was run on VM~270. The phase establishes
that the reference machine was unsuitable for direct reproduction, not which
one of its two protective conditions would independently have stopped a trial.
\end{evidencebox}

\evidencepair{01-vulnerability-validation/screenshots/01-validation-integrity-check.png}{01-vulnerability-validation/screenshots/02-patched-kernel-and-kmod.png}{Phase 01 reference-validation screenshots, visual page 1 of 4.}
\evidencepair{01-vulnerability-validation/screenshots/03-algif-aead-mitigation-active.png}{01-vulnerability-validation/screenshots/04-algif-aead-module-and-kernel-config.png}{Phase 01 reference-validation screenshots, visual page 2 of 4.}
\evidencepair{01-vulnerability-validation/screenshots/05-algif-aead-not-loaded.png}{01-vulnerability-validation/screenshots/06-local-transfer-integrity-check.png}{Phase 01 reference-validation screenshots, visual page 3 of 4.}
\evidencesingle{01-vulnerability-validation/screenshots/07-screenshot-integrity-check.png}{Phase 01 reference-validation screenshots, visual page 4 of 4.}
\clearpage

\subsection{Phase 02: Controlled Reproduction}
\label{app:evidence-phase02}

\begin{evidencebox}
\textbf{Public path:} \path{evidence/02-reproduction/}

\textbf{Manifest:} \path{evidence/02-reproduction/PUBLIC-SHA256SUMS}

\textbf{Coverage:} five machine-readable artifacts and six screenshots
covering the clean pre-state, module state, UID~0 result, post-shell state,
mounted-versus-backing view, and reboot recovery. Screenshot ID~04 remains
intentionally absent because it exposed the complete proof-of-concept source.

\textbf{Boundary:} the 11 captured objects verify; the README's listed digest
does not. The preserved vulnerable launch has no immediate pre-launch sample
hash, so exact Python-source identity is documentation-level provenance. The
different target digests establish observed read-path divergence, not a complete
block-forensic or on-disk replacement claim. Reboot recovery does not remediate
the vulnerable kernel.
\end{evidencebox}

\evidencepair{02-reproduction/screenshots/01-pre-exploit-state.png}{02-reproduction/screenshots/02-algif-aead-loaded.png}{Phase 02 controlled-reproduction screenshots, visual page 1 of 3.}
\evidencepair{02-reproduction/screenshots/03-privilege-escalation-success.png}{02-reproduction/screenshots/05-post-exploit-state.png}{Phase 02 controlled-reproduction screenshots, visual page 2 of 3.}
\evidencepair{02-reproduction/screenshots/06-page-cache-vs-backing-filesystem.png}{02-reproduction/screenshots/07-reboot-restores-original-file.png}{Phase 02 controlled-reproduction screenshots, visual page 3 of 3.}
\clearpage

\subsection{Phase 03: Detection and Exact-Sample Comparison}
\label{app:evidence-phase03}

\begin{evidencebox}
\textbf{Public path:} \path{evidence/03-detection/}

\textbf{Comparison path:}
\path{evidence/03-detection/yara-comparison/}

\textbf{Coverage:} detector artifacts, the MORI implementation and systemd
unit, 13 detector screenshots, five YARA artifacts, and two YARA screenshots.
All 28 captured objects verify; only the two listed README digests fail.

\textbf{Boundary:} MORI establishes a protected-target integrity deviation and
the exact-sample comparison establishes one match/non-match pair. Neither alone
attributes every deviation to Copy Fail, measures detection latency, or
evaluates YARA beyond the exact-hash rule used here.
\end{evidencebox}

\evidencepair{03-detection/screenshots/00-suid-target-surface.png}{03-detection/screenshots/01-suid-accessibility.png}{Phase 03 MORI detection screenshots, visual page 1 of 7.}
\evidencepair{03-detection/screenshots/02-suid-known-good-baseline.png}{03-detection/screenshots/03-fusermount3-usrmerge-verification.png}{Phase 03 MORI detection screenshots, visual page 2 of 7.}
\evidencepair{03-detection/screenshots/04-protected-baseline.png}{03-detection/screenshots/05-baseline-tamper-denied.png}{Phase 03 MORI detection screenshots, visual page 3 of 7.}
\evidencesidebyside{03-detection/screenshots/06a-mori-integrity-watcher-poc-part1.png}{03-detection/screenshots/06b-mori-integrity-watcher-poc-part2.png}{Phase 03 MORI detection screenshots, visual page 4 of 7.}
\evidencepair{03-detection/screenshots/07-mori-integrity-monitor-baseline.png}{03-detection/screenshots/09-mori-systemd-service-active.png}{Phase 03 MORI detection screenshots, visual page 5 of 7.}
\evidencepair{03-detection/screenshots/10-mori-system-service-clean-state.png}{03-detection/screenshots/11-mori-detects-copyfail.png}{Phase 03 MORI detection screenshots, visual page 6 of 7.}
\evidencesingle{03-detection/screenshots/12-mori-detection-with-root-confirmation.png}{Phase 03 MORI detection screenshots, visual page 7 of 7.}
\evidencepair{03-detection/yara-comparison/screenshots/00-yara-known-poc-match.png}{03-detection/yara-comparison/screenshots/01-yara-harmless-variant-no-match.png}{Phase 03 exact-sample YARA comparison, visual page 1 of 1.}
\clearpage

\subsection{Phase 04: Custom Mitigation and MORI v2.7}
\label{app:evidence-phase04}

\begin{evidencebox}
\textbf{Public path:} \path{evidence/04-custom-mitigation/}

\textbf{Manifest:}
\path{evidence/04-custom-mitigation/PUBLIC-SHA256SUMS.txt}

\textbf{Coverage:} 126 verified objects: ten MORI Guard v1 artifacts, eight v1
screenshots, 49 MORI v2 implementation or provenance artifacts, and 59 MORI v2
screenshots. The screenshot sequence retains false positives, negative
controls, historical checkpoints, the v2.6.2 expiry bypass, the first v2.7
failure, the corrected regression, two final PoC paths, and final provenance.

\textbf{Boundary:} the frozen v2.7 result applies to the tested same-process
paths and controls. A later provenance capture on the patched runtime verifies
the transferred artifact hashes; it does not replace the vulnerable-kernel test
context or establish universal prevention and compatibility.
\end{evidencebox}

\subsubsection{MORI Guard v1 visual index}

\evidencepair{04-custom-mitigation/screenshots/00-mori-guard-configuration.png}{04-custom-mitigation/screenshots/01-mori-guard-denies-module-load.png}{Phase 04 MORI Guard v1 screenshots, visual page 1 of 4.}
\evidencepair{04-custom-mitigation/screenshots/02-mori-guard-final-message.png}{04-custom-mitigation/screenshots/03-pre-mitigation-test-state.png}{Phase 04 MORI Guard v1 screenshots, visual page 2 of 4.}
\evidencepair{04-custom-mitigation/screenshots/04-copyfail-blocked-by-mori-guard.png}{04-custom-mitigation/screenshots/05-mori-guard-block-event.png}{Phase 04 MORI Guard v1 screenshots, visual page 3 of 4.}
\evidencepair{04-custom-mitigation/screenshots/06-mori-guard-post-attempt-verification.png}{04-custom-mitigation/screenshots/07-mori-guard-blocker-source.png}{Phase 04 MORI Guard v1 screenshots, visual page 4 of 4.}
\clearpage

\subsubsection{MORI v2 development and final-validation visual index}

\evidencepair{04-custom-mitigation/mori-v2/screenshots/00-mori-v2-identifies-mori-monitor-process.png}{04-custom-mitigation/mori-v2/screenshots/01-mori-v2-false-positive-mori-monitor.png}{Phase 04 MORI v2 screenshots, visual page 1 of 30.}
\evidencepair{04-custom-mitigation/mori-v2/screenshots/02-mori-monitor-service-context.png}{04-custom-mitigation/mori-v2/screenshots/03-mori-v2-1-1-trusted-monitor-loaded.png}{Phase 04 MORI v2 screenshots, visual page 2 of 30.}
\evidencepair{04-custom-mitigation/mori-v2/screenshots/04-mori-v2-1-1-trusted-monitor-suppressed.png}{04-custom-mitigation/mori-v2/screenshots/05-mori-v2-1-1-untrusted-suid-read-detected.png}{Phase 04 MORI v2 screenshots, visual page 3 of 30.}
\evidencepair{04-custom-mitigation/mori-v2/screenshots/06-mori-v2-1-1-implementation-hashes.png}{04-custom-mitigation/mori-v2/screenshots/07-mori-v2-1-2-monitor-trust-verified.png}{Phase 04 MORI v2 screenshots, visual page 4 of 30.}
\evidencepair{04-custom-mitigation/mori-v2/screenshots/08-mori-v2-1-2-verified-monitor-suppressed.png}{04-custom-mitigation/mori-v2/screenshots/09-mori-v2-1-2-untrusted-suid-read-detected.png}{Phase 04 MORI v2 screenshots, visual page 5 of 30.}
\evidencepair{04-custom-mitigation/mori-v2/screenshots/10-mori-v2-1-2-tampered-monitor-trust-refused.png}{04-custom-mitigation/mori-v2/screenshots/11-mori-v2-1-2-refused-monitor-becomes-visible.png}{Phase 04 MORI v2 screenshots, visual page 6 of 30.}
\evidencepair{04-custom-mitigation/mori-v2/screenshots/12-mori-v2-1-2-restored-monitor-trust-recovered.png}{04-custom-mitigation/mori-v2/screenshots/13-mori-v2-1-2-final-artifact-hashes.png}{Phase 04 MORI v2 screenshots, visual page 7 of 30.}
\evidencepair{04-custom-mitigation/mori-v2/screenshots/14-mori-v2-2-aead-bind-attempt-observed.png}{04-custom-mitigation/mori-v2/screenshots/15-mori-v2-2-v1-guard-intercepts-algif-aead.png}{Phase 04 MORI v2 screenshots, visual page 8 of 30.}
\evidencepair{04-custom-mitigation/mori-v2/screenshots/16-mori-v2-2-benign-aead-bind-preserved.png}{04-custom-mitigation/mori-v2/screenshots/17-mori-v2-2-final-artifact-hashes.png}{Phase 04 MORI v2 screenshots, visual page 9 of 30.}
\evidencepair{04-custom-mitigation/mori-v2/screenshots/18-mori-v2-3-ktime-helper-verifier-probe.png}{04-custom-mitigation/mori-v2/screenshots/19-mori-v2-3-same-tgid-aead-suid-correlation.png}{Phase 04 MORI v2 screenshots, visual page 10 of 30.}
\evidencepair{04-custom-mitigation/mori-v2/screenshots/20-mori-v2-3-correlation-window-expires.png}{04-custom-mitigation/mori-v2/screenshots/21-mori-v2-3-suid-only-negative-control.png}{Phase 04 MORI v2 screenshots, visual page 11 of 30.}
\evidencepair{04-custom-mitigation/mori-v2/screenshots/22-mori-v2-3-final-artifact-hashes.png}{04-custom-mitigation/mori-v2/screenshots/23-mori-v2-4-benign-file-pipe-splice-observed.png}{Phase 04 MORI v2 screenshots, visual page 12 of 30.}
\evidencepair{04-custom-mitigation/mori-v2/screenshots/24-mori-v2-4-suid-file-pipe-splice-observed.png}{04-custom-mitigation/mori-v2/screenshots/25-mori-v2-4-splice-probe-hashes.png}{Phase 04 MORI v2 screenshots, visual page 13 of 30.}
\evidencepair{04-custom-mitigation/mori-v2/screenshots/26-mori-v2-4-cross-hook-splice-state-correlation.png}{04-custom-mitigation/mori-v2/screenshots/27-mori-v2-4-aead-splice-suid-triple-correlation.png}{Phase 04 MORI v2 screenshots, visual page 14 of 30.}
\evidencepair{04-custom-mitigation/mori-v2/screenshots/28-mori-v2-4-final-artifact-hashes.png}{04-custom-mitigation/mori-v2/screenshots/29-mori-v2-5-shadow-policy-would-deny.png}{Phase 04 MORI v2 screenshots, visual page 15 of 30.}
\evidencepair{04-custom-mitigation/mori-v2/screenshots/30-mori-v2-5-aead-only-negative-control.png}{04-custom-mitigation/mori-v2/screenshots/31-mori-v2-5-splice-only-negative-control.png}{Phase 04 MORI v2 screenshots, visual page 16 of 30.}
\evidencepair{04-custom-mitigation/mori-v2/screenshots/32-mori-v2-5-final-artifact-hashes.png}{04-custom-mitigation/mori-v2/screenshots/33-mori-v2-6-selective-enforcement-blocks-correlated-splice.png}{Phase 04 MORI v2 screenshots, visual page 17 of 30.}
\evidencepair{04-custom-mitigation/mori-v2/screenshots/34-mori-v2-6-final-artifact-hashes.png}{04-custom-mitigation/mori-v2/screenshots/35-mori-v2-6-1-ring-buffer-deny-event-delivery.png}{Phase 04 MORI v2 screenshots, visual page 18 of 30.}
\evidencepair{04-custom-mitigation/mori-v2/screenshots/36-mori-v2-6-1-final-artifact-hashes.png}{04-custom-mitigation/mori-v2/screenshots/37a-mori-v2-6-2-attacker-facing-judgment-cat.png}{Phase 04 MORI v2 screenshots, visual page 19 of 30.}
\evidencepair{04-custom-mitigation/mori-v2/screenshots/37b-mori-v2-6-2-attacker-facing-judgment-cat.png}{04-custom-mitigation/mori-v2/screenshots/38-mori-v2-6-2-final-artifact-hashes.png}{Phase 04 MORI v2 screenshots, visual page 20 of 30.}
\evidencepair{04-custom-mitigation/mori-v2/screenshots/39-mori-v2-6-2-pre-poc-validation.png}{04-custom-mitigation/mori-v2/screenshots/40-mori-v2-6-2-active-before-real-poc.png}{Phase 04 MORI v2 screenshots, visual page 21 of 30.}
\evidencepair{04-custom-mitigation/mori-v2/screenshots/41-mori-v2-6-2-real-copy-fail-poc-blocked.png}{04-custom-mitigation/mori-v2/screenshots/42-mori-v2-6-2-post-poc-validation-clean.png}{Phase 04 MORI v2 screenshots, visual page 22 of 30.}
\evidencepair{04-custom-mitigation/mori-v2/screenshots/43-mori-v2-6-2-c-poc-pre-state.png}{04-custom-mitigation/mori-v2/screenshots/44-mori-v2-6-2-c-poc-denied.png}{Phase 04 MORI v2 screenshots, visual page 23 of 30.}
\evidencepair{04-custom-mitigation/mori-v2/screenshots/45-mori-v2-6-2-c-poc-post-validation.png}{04-custom-mitigation/mori-v2/screenshots/46-alternate-suid-target-baseline-fusermount3.png}{Phase 04 MORI v2 screenshots, visual page 24 of 30.}
\evidencepair{04-custom-mitigation/mori-v2/screenshots/47-mori-v2-6-2-suid-target-sweep-denied.png}{04-custom-mitigation/mori-v2/screenshots/48-mori-v2-6-2-post-suid-sweep-integrity-10-of-10.png}{Phase 04 MORI v2 screenshots, visual page 25 of 30.}
\evidencepair{04-custom-mitigation/mori-v2/screenshots/49-mori-v2-6-2-cross-tgid-isolation-allowed.png}{04-custom-mitigation/mori-v2/screenshots/50-mori-v2-6-2-same-tgid-expiry-bypass-reproduced.png}{Phase 04 MORI v2 screenshots, visual page 26 of 30.}
\evidencepair{04-custom-mitigation/mori-v2/screenshots/51-mori-v2-7-lifecycle-regression-splice-still-allowed.png}{04-custom-mitigation/mori-v2/screenshots/52-mori-v2-7-expiry-bypass-regression-fixed.png}{Phase 04 MORI v2 screenshots, visual page 27 of 30.}
\evidencepair{04-custom-mitigation/mori-v2/screenshots/53-mori-v2-7-final-python-poc-denied.png}{04-custom-mitigation/mori-v2/screenshots/54-mori-v2-7-final-python-post-validation.png}{Phase 04 MORI v2 screenshots, visual page 28 of 30.}
\evidencepair{04-custom-mitigation/mori-v2/screenshots/55-mori-v2-7-final-c-poc-denied-and-post-validation.png}{04-custom-mitigation/mori-v2/screenshots/56-independent-c-poc-provenance.png}{Phase 04 MORI v2 screenshots, visual page 29 of 30.}
\evidencesingle{04-custom-mitigation/mori-v2/screenshots/57-mori-v2-7-final-artifact-provenance.png}{Phase 04 MORI v2 screenshots, visual page 30 of 30.}
\clearpage

\subsection{Phase 05: Vendor-Patched Validation}
\label{app:evidence-phase05}

\begin{evidencebox}
\textbf{Public path:} \path{evidence/05-vendor-patch-validation/}

\textbf{Manifest:}
\path{evidence/05-vendor-patch-validation/PUBLIC-SHA256SUMS.txt}

\textbf{Coverage:} five verified text artifacts and 13 verified screenshots
covering the patched baseline, MORI isolation, kernel package, benign AF\_ALG
control, Python retest, independent C provenance and retest, and final target
validation.

\textbf{Boundary:} the pre-patch filename is inaccurate; C hashes come from
screenshot~10 because the text capture failed; strict isolation-before-C-test
order is not proven by timestamps alone; the patched Python account differs
from the vulnerable run; and the interactive retests have no matching raw
execution transcripts. The benign control proves socket creation, bind, and
accept availability, not encryption or decryption. The independent C path has
no preserved successful unmitigated vulnerable-side run, and the final target
check establishes state at capture time rather than continuous integrity.
\end{evidencebox}

\evidencepair{05-vendor-patch-validation/screenshots/00-vendor-patch-baseline.png}{05-vendor-patch-validation/screenshots/01-vendor-patch-baseline-hash.png}{Phase 05 vendor-patched validation screenshots, visual page 1 of 7.}
\evidencepair{05-vendor-patch-validation/screenshots/02-mori-enforcement-isolated.png}{05-vendor-patch-validation/screenshots/03-vendor-kernel-package-version.png}{Phase 05 vendor-patched validation screenshots, visual page 2 of 7.}
\evidencepair{05-vendor-patch-validation/screenshots/04-algif-aead-module-state.png}{05-vendor-patch-validation/screenshots/05-benign-aead-control-pass.png}{Phase 05 vendor-patched validation screenshots, visual page 3 of 7.}
\evidencepair{05-vendor-patch-validation/screenshots/06-benign-aead-control-hash.png}{05-vendor-patch-validation/screenshots/07-python-poc-provenance-and-denial.png}{Phase 05 vendor-patched validation screenshots, visual page 4 of 7.}
\evidencepair{05-vendor-patch-validation/screenshots/08-post-python-target-integrity.png}{05-vendor-patch-validation/screenshots/09-pre-c-poc-target-integrity.png}{Phase 05 vendor-patched validation screenshots, visual page 5 of 7.}
\evidencepair{05-vendor-patch-validation/screenshots/10-independent-c-poc-provenance.png}{05-vendor-patch-validation/screenshots/11-independent-c-poc-denied.png}{Phase 05 vendor-patched validation screenshots, visual page 6 of 7.}
\evidencesingle{05-vendor-patch-validation/screenshots/12-post-c-poc-integrity-auth-control.png}{Phase 05 vendor-patched validation screenshots, visual page 7 of 7.}
\clearpage

\subsection{Cross-Chapter Consistency Result}
\label{app:consistency-result}

The final source was checked across Chapters~1--10 for virtual-machine roles,
phase numbering, kernel and package versions, user identities, target and
sample hashes, MORI versions, evidence counts, research-question answers, and
the Phase~05 caveats. \Cref{tab:cross-chapter-consistency} records the
authoritative values used after the audit.

\begin{table}[H]
    \centering
    \caption{Cross-chapter values fixed by the evidence audit.}
    \label{tab:cross-chapter-consistency}
    \small
    \begin{tabularx}{\textwidth}{@{}>{\bfseries\raggedright\arraybackslash}p{0.27\textwidth}Y@{}}
        \toprule
        Item & Authoritative report value \\
        \midrule
        Reference / experimental VMs & VM~270 \code{copyfail01} is the patched reference; VM~271 \code{copyfail-vuln} is the state-changing target. \\
        Vulnerable runtime & \code{6.8.0-116-generic}, package \code{6.8.0-116.116}. \\
        First fixed Noble generic package & \code{6.8.0-117.117}, as identified by Canonical. \\
        Final patched runtime & \code{6.8.0-137-generic}, package \code{6.8.0-137.137}. \\
        Test identities & Vulnerable and MORI tests use \code{labuser}, UID~1001; the patched Python screenshot uses \code{researcher}, UID~1000. \\
        Known-good target digest & {\scriptsize\ttfamily c74311fe5636b7d7f9a56239fa8adeeab12ba86fe7d41b91afa85bf9bbdae78b}. \\
        Changed mounted-view digest & {\scriptsize\ttfamily 44900c631391f0d60eb6d271b8374a08dc1d9be76e403390d27a91ed5f179be9}. \\
        Documented Python sample & Phase~02 documentation attributes 731 bytes and SHA-256 {\scriptsize\ttfamily d401e7d1c00605749d6c617ace73ab20a762b72e41c2e1590331596e38219a61} to the sample associated with the UID~0 result, but the vulnerable launch has no immediate pre-launch hash. Phase~05 captures that digest before its patched-side execution. The non-executed trailing-newline variant is 732 bytes. \\
        Frozen MORI release & v2.7 executable SHA-256 {\scriptsize\ttfamily 7fe7b609b90e164f3ae4a9c025363399a8a63bf2a4338db2ecc62cfd2682d039}. \\
        MORI v2 inventory & 108 verified objects: 49 implementation/provenance artifacts and 59 screenshots. \\
        Patched validation inventory & 18 verified objects: five text artifacts and 13 screenshots. \\
        Outcome boundary & Demonstrated reproduction and tested-path prevention are empirical laboratory results. MORI does not provide universal prevention, and the patched retest is not a population-wide statistical result. The C path is additional patched-side coverage, not a matched vulnerable/patched differential. \\
        \bottomrule
    \end{tabularx}
\end{table}

The audit does not remove inconvenient records. The missing reproduction and
detection IDs, stale documentation hashes, v2.6.2 bypass, first v2.7 regression
failure, later provenance environment, and Phase~05 chronology defects remain
visible because they determine the strength of the conclusions. The assembled
package includes a separate claim-to-evidence map and audit README so the same
decisions can be checked without reverse-engineering them from the prose.